\section{Introduction}

Secret-key distillation asks whether two distant laboratories can turn a
shared noisy resource into identical random bits that an eavesdropper cannot
predict.  For a bipartite quantum state $\rho_{AB}$, Alice and Bob act locally
and exchange authenticated public classical messages.  Eve holds a
purification of the input and receives the complete public transcript.  An
ideal key of length $\ell$ is uniform, agrees at the two laboratories, and is
independent of Eve: its joint state has the form
\[
 2^{-\ell}\sum_{z\in\{0,1\}^{\ell}}
 |z,z\rangle\langle z,z|\otimes\sigma_{E'}.
\]
We use the usual ideal-key security criterion
\cite{keyPrivateStates,RennerKonig05}.  The input is the only shared private
resource.  Alice and Bob may process all copies jointly in their respective
laboratories and use local auxiliaries and private randomness.
They may also begin with a public random variable, independent of the input
and its reference and recorded for Eve.  Conditional on that public
variable, the input-reference state and the two private auxiliary resources
form a product, with fixed preparation kernels for every input.

An \emph{admitted} protocol has a finite or $\mathbb N$-indexed sequence of
measurable local-instrument steps, after any supplied local aggregation.
Its histories and terminal readouts must satisfy the reference-compatibility
conditions of Definition~\ref{def:protocol}.  These include ordinary finite
protocols with finite classical records obeying the preceding resource rules,
as well as the specified countable processes.  Eve retains the initial
purification and the full public record, including speaker choices and
stopping information.  A protocol is \emph{completed on the actual input}
when both laboratories produce their prescribed outputs almost surely and
the joint output is normalized.

The distillable secret key $K_D(\rho)$ is the supremum of
$\liminf_n\ell_n/n$ over admitted completed key-output sequences on
$\rho^{\otimes n}$ whose error from an ideal key tends to zero
(Definition~\ref{def:distillable-key}).  All errors use the unhalved trace
norm $\norm{X}_1=\tr\sqrt{X^*X}$.

We construct an entangled state for which even one approximately secret bit
is uniformly out of reach.  A density operator is a positive matrix of trace
one.  It is \emph{separable} if it is a finite sum of positive product
matrices, and \emph{entangled} otherwise.

\begin{theorem}\label{thm:main-state}
There is an explicitly specified entangled density operator $\rho$ on
$\C^{10}\otimes\C^{10}$ whose range contains no nonzero product vector and
for which $K_D(\rho)=0$.  More precisely, let $n\geq1$ and let
$\tau_{K_AK_BE'}$ be the normalized output of any admitted completed bit-output protocol of
Definition~\ref{def:protocol} on $\rho^{\otimes n}$, where $E'$ retains an
initial purification and the complete public transcript.  For every state
$\sigma_{E'}$,
\begin{equation}\label{eq:main-gap}
 \left\lVert \tau_{K_AK_BE'}-
 \frac12\sum_{i=0}^1|i,i\rangle\langle i,i|\otimes\sigma_{E'}
 \right\rVert_1\geq\frac15.
\end{equation}
Thus the corresponding trace distance, with its conventional factor
$1/2$, is at least $1/10$.
\end{theorem}

There is no uniform bound on the amount of communication or the number of
steps within the admitted model.  Proposition~\ref{prop:operational-extensions}
also gives separate consequences for public conditioning on an event of
positive probability, vanishing abort probability, and trace-norm limits of
completed outputs whose Eve registers retain or allow recovery of the
purification and complete public record.

The state comes from a construction for completely positive maps.  A
complex-linear map $F:\Mat a\to\Mat b$ is \emph{completely positive} (CP)
when $\mathrm{id}_k\otimes F$ preserves positive matrices for every
$k\geq1$.  We call it \emph{PPT} when both $F$ and $T_b\circ F$ are CP,
where $T_b$ is coordinate transposition.  Write $J(F)$ for the unnormalized
input-first Choi matrix, defined in \eqref{eq:choi-definition}.  A CP map is
\emph{entanglement breaking} when applying it to one subsystem of any
positive bipartite input always gives a separable output.  Equivalently, its
Choi matrix is separable \cite[Theorem 4]{HSR2003}; this cone formulation
does not require trace preservation.

\begin{theorem}\label{thm:main-pair}
There are explicitly specified PPT maps
$\Phi_1,\Phi_2:\Mat{10}\to\Mat{10}$ for which
\[
 Z=J(\Phi_2\circ\Phi_1)
\]
is nonzero and its range contains no nonzero product vector in
$\C^{10}\otimes\C^{10}$.  In particular, the composition
$\Phi_2\circ\Phi_1$ is not entanglement breaking.
\end{theorem}

The maps are given by \eqref{eq:phi-pair}, starting from the integer matrices
in \eqref{eq:explicit-matrices}.  They may differ, and the construction does
not impose trace preservation or unitality.  The state in
Theorem~\ref{thm:main-state} is
\begin{equation}\label{eq:main-state-formula}
                    \rho=\frac{Z}{\tr Z}.
\end{equation}
The denominator is positive because $Z$ is a nonzero positive Choi matrix.
A separate construction supplies a trace-preserving example.

\begin{theorem}\label{thm:main-channel}
There is an explicitly specified trace-preserving PPT channel
$\Theta:\Mat{21}\to\Mat{21}$ whose square $\Theta\circ\Theta$ is not
entanglement breaking.
\end{theorem}

The pair in Theorem~\ref{thm:main-pair} gives a negative answer to the
unrestricted two-map PPT-composition conjecture of Christandl,
M\"uller-Hermes, and Wolf \cite[Conjecture IV.1]{CMHW2019}.  The channel in
Theorem~\ref{thm:main-channel} answers the channel question recorded in the
2012 BIRS workshop report \cite[Problem G]{BIRS2012}, already with the same
channel in both positions.  We do not claim that either dimension is
minimal.

\paragraph{Entanglement and secrecy.}
Public discussion and privacy amplification are central to the classical
theory of key agreement \cite{BennettBrassardRobert88,Maurer93,BBCM95}.
Entanglement purification provides a quantum route to secrecy: nearly pure
shared entangled pairs decouple from a purifying adversary
\cite{BBPSSW96,BDSW96,DeutschEtAl96}.  Bound entanglement showed that
entanglement need not be distillable into singlets.  In particular, PPT
entangled states have zero distillable entanglement \cite{HHH98}.  Some PPT
entangled states nevertheless have positive distillable secret key
\cite{keyBoundEntanglement}.  Their private-state description explains how
secrecy can persist in a quantum shield even when singlet distillation is
impossible \cite{keyPrivateStates}.

Separable states have zero key against a purifying Eve
\cite{CurtyLewensteinLutkenhaus04,keyBoundEntanglement}.  The converse was
posed explicitly as Problem 24, ``Secret key from all entangled states,'' in
Kr\"uger and Werner's open-problem collection: the entry is dated
15 March 2005 and lists P.~Horodecki as contact \cite{KrugerWerner05}.
The question is also discussed in the private-state literature
\cite{keyPrivateStates,HorodeckiEtAl18}.  One-way coding theorems give useful
positive-key criteria \cite{DevetakWinter05}, while Theorem~\ref{thm:main-state}
in particular rules out a positive key rate for a specific entangled state
under ordinary finite two-way protocols, whose round count may grow with
the input length.  The distinction
between secrecy needed to prepare correlations and key recoverable from
them also matters: measurements of an entangled state can yield correlations
with a positive secrecy formation cost without proving a positive
distillation rate \cite{GisinWolf00,GisinRennerWolf02,AcinGisin05}.

There are earlier quantitative obstructions near private bits.  Kim and
Sanders obtained a distance bound for PPT states under a Hermitian-block
hypothesis \cite{KimSanders10}, and Badzi\k{a}g and coauthors removed that
hypothesis \cite{BadziagEtAl14}.  Those bounds depend on the shield dimension,
which may grow with a protocol.  Our bound applies to the final classical bit
registers and Eve's complete public side information, with a constant
independent of the input length and the local memories.

Pauwels, Gisin, and Renner recently constructed a finite classical bipartite
source with positive secrecy formation cost and zero asymptotic secret-key
rate under finite-round authenticated public discussion, with the round
count allowed to grow with block length
\cite[Theorem 1 and Appendix A]{PauwelsGisinRenner2026}.
Their revised paper also proves positive two-way distillable quantum key for
the Horodecki two-qutrit states $\sigma_\alpha$ with $3<\alpha\leq5$,
including the bound-entangled range $3<\alpha\leq4$, against a purifying quantum Eve
\cite[Corollary 3 and Appendix I.7]{PauwelsGisinRenner2026}.  The state in
Theorem~\ref{thm:main-state} gives the quantum zero-key conclusion under the
local quantum protocol model above.

\paragraph{The operational obstruction.}
The first part of the proof concerns a class $\mathscr C$ of states
$\omega$ on $\C^a\otimes\C^b$ with a local representation
\[
 \omega=(\mathcal L\otimes\mathcal M)(\chi\chi^*),
 \quad
 \chi\in\C^r\otimes\C^s,\quad
 \mathcal L:\Mat r\to\Mat a,\quad
 \mathcal M:\Mat s\to\Mat b,
\]
where $\mathcal L,\mathcal L\circ T_r,\mathcal M,\mathcal M\circ T_s$
are CP.  These input-transpose conditions are equivalent to the
output-transpose PPT convention above, as Lemma~\ref{lem:ppt-stability}
shows.  The maps may be rectangular and need not preserve trace; only
the output $\omega$ has trace one.  This condition persists under tensor powers.
Relative to one common transcript measure, the four final bit branches have
positive local effect densities $X_i(t)\otimes Y_j(t)$ on the two local
input spaces of the protocol for almost every
complete transcript $t$, with the same two local pairs for all
$i,j\in\{0,1\}$.  To prove this, we run the protocol on auxiliary
product probes.  Their two local sampling tapes are conditionally independent
given the complete public record.  Informationally complete tests on
untouched references turn that independence into common product factors,
and a filter on the references transfers the factors to the actual input.

For a state in $\mathscr C$, a coherence inequality bounds the overlap
between the $00$ and $11$ branches by the probabilities of disagreement.
After integration over transcripts, let $e$ be the total probability of
different bits and let $\mathfrak f$ be the root fidelity between Eve's
unnormalized $00$ and $11$ conditional states, summed over the public
record.  We prove
\[
                  \mathfrak f\leq\sqrt{e(1-e)}+e.
\]
If the output is at trace-norm distance $\eta$ from an ideal bit, correctness
and secrecy give, respectively,
\[
                 e\leq\frac\eta2,
       \qquad \mathfrak f\geq\frac{1-\eta}{2}.
\]
For $\eta<1/5$, the right side of the first display is less than $2/5$,
whereas the second display requires more than $2/5$.  This proves the uniform
gap.  Taking the first bit of a longer ideal key is a trace-norm contraction,
so no positive asymptotic key rate is possible.

\paragraph{How the dimension-ten construction enters.}
For CP maps $F:\Mat a\to\Mat b$ and $G:\Mat b\to\Mat c$, define
$F^\sharp=T_a\circ F^\dagger\circ T_b$, where $F^\dagger$ is the
Hilbert--Schmidt adjoint.  The transpose-adjoint Choi identity is
\[
 J(G\circ F)=(F^\sharp\otimes G)(\Omega_b\Omega_b^*),
 \qquad \Omega_b=\sum_{i=1}^b e_i\otimes e_i.
\]
This identity and the local-product formulation are established machinery
\cite[Lemma I.1(2), Section IV.B]{CMHW2019}; we prove the needed rectangular
form and normalization in Lemma~\ref{lem:choi-transfer}.  When $F$ and $G$
are PPT and the Choi matrix is nonzero, division by its trace gives a state
in $\mathscr C$.  Applying this observation to
Theorem~\ref{thm:main-pair} proves that \eqref{eq:main-state-formula} belongs
to the class covered by the operational obstruction.

It remains to make this state entangled.  A nonzero separable positive
matrix has a nonzero product vector in its range, a necessary part of the
range criterion \cite{Horodecki1997}.  We exclude such vectors directly.
Start with a PPT map $L:\Mat4\to\Mat4$ and twenty projective directions
$[x]$ at which $L(xx^*)$ is singular.  No nonzero homogeneous quadratic in
four variables vanishes at ten of these directions.  The ten-dimensional
symmetric square records the quadratic monomials of $x$, and the
six-dimensional exterior square converts rank loss into a product vector
in $\ker Z$.  If $\ran Z$ contained a product vector, its two factors would
give two nonzero quadratics whose zero sets cover all twenty directions.
Each can cover at most nine, and $9+9<20$.  A separate positive definite
output of $L$ ensures that $Z$ is nonzero.

This range argument is related to unextendible product bases
\cite{Bennett1999}, nonorthogonal product exclusion by linear independence
\cite[Section IV, Lemma 3]{Pittenger2003}, and polynomial evaluation methods
\cite[Proposition 1.3]{Parthasarathy2004}.  The complementary-pair operation
on the exterior square agrees with the four-dimensional fermionic duality
of Schliemann, Loss, and MacDonald
\cite[Section III]{SchliemannLossMacDonald2001}.  Our explicit $L$ comes from
a real integer $6$-by-$4$ matrix pencil.  Its rank-losing directions solve a
bilinear system of the type appearing in the study of low-rank PPT states
on $\C^4\otimes\C^4$ \cite[Section 7]{HansenEtAl2012}.  Exact reductions at
three primes prove that the relevant degree-twenty algebra is a field with
normal-closure Galois group $S_{20}$.  This symmetry carries one nonzero
quadratic evaluation determinant to every ten-point choice.  The arithmetic
certificate and its precise computational coverage are given in
Appendix~\ref{app:arithmetic}.

\paragraph{The channel and the projection construction.}
The PPT-composition problem has important affirmative cases, including two
maps on $\Mat3$ \cite[Corollary III.1]{CMHW2019}; see also
\cite{ChenYangTang2019}.  Other cases include Gaussian channels
\cite{CMHW2019} and diagonal unitary covariant or conjugate diagonal
unitary covariant PPT maps \cite[Theorem~4.5]{SinghNechita2022}.
Asymptotic approach to the entanglement-breaking cone was proved for
iterates of unital or trace-preserving PPT maps
\cite[Theorem~3.3]{KMP2018}.  Finite entanglement-breaking index was
proved for unital PPT channels by Rahaman, Jaques, and Paulsen
\cite[Theorem~4.4]{RahamanJaquesPaulsen2018}, and for PPT channels with a
full-rank invariant state by Hanson, Rouz\'e, and Stilck Fran\c ca
\cite[Theorem~3.10]{HansonRouzeFranca2020}.
More recent results concern different scopes.  An and Lee prove that a
qutrit CP map whose Choi matrix is one-copy undistillable and a qutrit CP
map whose Choi matrix has Schmidt number at most two compose to an
entanglement-breaking map in either order \cite[Theorem~4]{AnLee2026}.
Park's preprint states eventual entanglement breaking for each CP and
completely copositive endomap, with a map-dependent number of iterations
\cite[Theorem~1.1]{Park2026}.  A failure at the second iterate is
compatible with such eventual behavior.

To obtain Theorem~\ref{thm:main-channel}, we rescale the two maps, alternate
them between two coordinate blocks, and send the missing trace to an
absorbing flag.  The switch follows Christandl, M\"uller-Hermes, and Wolf
\cite[Section IV.A]{CMHW2019}; the flag completion follows Filippov's
trace-preserving extension \cite[Proposition 3, Equation (15)]{Filippov2021}.
The Choi matrix of the square has a local corner equal to a positive multiple
of $Z$.  Local compression preserves separability, so that corner certifies
that the square is not entanglement breaking.

The final part of the paper develops a second way to retain entanglement
under two PPT maps.  A rectangular criterion starts from real unit vectors
and projective measurements, uses a block Gram matrix for positivity, and
uses commutation for the partial-transpose property.  A support argument
using zero tensor-basis diagonal entries rules out separability over the
complex field.  We realize
the criterion by a restricted-intersection construction
\cite{FW1981,ABS1991} and tensor products of the Mermin--Peres parity
measurements \cite{Mermin1990,Peres1990,Peres1991}.  Its dimensions are much
larger, but the mechanism is independent of the twenty-direction geometry.
Frankl and Rödl's forbidden-intersection theory similarly forces fixed-size
orthogonal families in large subsets of a sign cube
\cite[Theorem 1.11]{FR1987}; our eight-point estimate is proved by a direct
polynomial bound and an elementary tensor-power fiber argument.
The real construction also gives an entangled output when the same
trace-preserving rectangular PPT channel is applied to both halves of a
maximally entangled state.  This distinction between real matrices and complex
separability matters: the real-field example of Chiribella, Davidson,
Paulsen, and Rahaman has an entanglement-breaking complexification
\cite[Example 9.6]{CDPR2023}.

\paragraph{Organization.}
Section~\ref{sec:foundations} establishes the common map conventions,
represented class, and Choi transfer.  Sections~\ref{sec:coherence}--\ref{sec:secret-key}
prove the operational theorem.  Sections~\ref{sec:geometry} and
\ref{sec:certificate} construct the dimension-ten pair, and
Section~\ref{sec:applications} proves the state and channel consequences.
Sections~\ref{sec:projection} and \ref{sec:combinatorial} give the independent
projection and combinatorial method.  The appendices give the exact certificate, the general algebraic facts
used in its proof, and the supplementary finite checks of the explicit data.
